\section{Qué es un LLM Agent}

\subsection{Composición del sistema de un Agent}

\begin{figure}[H]
\centering
\includegraphics[page=4,width=0.9\textwidth]{slides.pdf}
\caption{El curso resume los componentes estándar de un Agent con Memory, Tool Use, Retrieval, Reasoning/Planning y Environment Feedback.}
\end{figure}

El panorama que presenta el curso es claro: el centro de un Agent sigue siendo el LLM, pero el LLM ya no funciona de manera aislada, sino que se acopla con sistemas externos mediante varias interfaces.

\begin{itemize}
    \item \textbf{Memory}: conserva el contexto de la tarea, las preferencias del usuario y las trayectorias históricas, para evitar que cada turno de conversación empiece desde cero.
    \item \textbf{Tool use}: permite que el modelo invoque búsquedas, ejecución de código, API, bases de datos u operaciones de navegador.
    \item \textbf{Retrieval}: extrae evidencia externa de documentos y bases de conocimiento, en lugar de depender por completo de la memoria almacenada en los parámetros.
    \item \textbf{Reasoning \& Planning}: divide objetivos complejos en varios pasos, en vez de entregar directamente una salida de un solo turno.
    \item \textbf{Feedback / Environment}: permite que el Agent revise sus propios resultados de ejecución, formando un ciclo cerrado de prueba y error.
\end{itemize}

\begin{knowledgebox}{Por qué conservar estos límites entre módulos}
Si la memoria, la recuperación, la invocación de herramientas y la planeación se comprimen todas en un prompt sumamente largo, el sistema será muy difícil de observar, depurar y optimizar. El valor principal de un framework de Agent está precisamente en separar de forma explícita estas capacidades, de modo que cada componente pueda sustituirse, restringirse y evaluarse por separado.
\end{knowledgebox}

\subsection{Tipos de entornos en los que actúa un Agent}

\begin{figure}[H]
\centering
\includegraphics[page=5,width=0.9\textwidth]{slides.pdf}
\caption{Un LLM Agent no existe solo en interfaces de chat, sino que también puede operar en páginas web, desarrollo de software, flujos de trabajo empresariales y robots, entre otros entornos diversos.}
\end{figure}

El curso hace especial énfasis en la diversidad de entornos de un Agent. Puede trabajar tanto en entornos digitales, como el IDE, el navegador o las plataformas de flujo de trabajo, como extenderse a entornos corpóreos, como robots y tareas de percepción multimodal. Cuanto más abierto es el entorno, más necesita el Agent:
\begin{itemize}
    \item tomar decisiones locales con información incompleta;
    \item recuperarse de los fallos, en vez de terminar de inmediato ante un solo error;
    \item coordinar al mismo tiempo operaciones simbólicas, interacción en lenguaje natural y acciones externas.
\end{itemize}

\subsection{Resumen de la sección}

La definición de un LLM Agent no es algo tan simple como ``un LLM capaz de invocar herramientas'', sino un sistema de ejecución de tareas compuesto por el modelo, capacidades externas y un ciclo de retroalimentación.

